\documentclass[UTF8,a4paper,11pt,fontset=fandol]{ctexart}
\usepackage[margin=2.2cm]{geometry}
\usepackage{booktabs,tabularx,array,enumitem,xcolor,fancyhdr,hyperref}
\definecolor{ink}{HTML}{244B66}
\hypersetup{colorlinks=true,linkcolor=ink,urlcolor=ink}
\pagestyle{fancy}\fancyhf{}\fancyhead[L]{CET-4 学习记录}\fancyhead[R]{2025-12-1}\fancyfoot[C]{\thepage}
\setlength{\headheight}{15pt}
\setlength{\parindent}{0pt}
\setlength{\parskip}{0.35em}
\linespread{1.12}
\setlist{nosep,leftmargin=1.7em}
\ctexset{section={format=\Large\bfseries\color{ink},beforeskip=1.0ex,afterskip=0.7ex},subsection={format=\normalsize\bfseries,beforeskip=0.8ex,afterskip=0.4ex}}
\newcommand{\en}[1]{\textit{#1}}
\begin{document}
\begin{center}
{\LARGE\bfseries 听力与作文练习复盘}\\[6pt]
{\large 听力第 1--2 题与作文初稿修改}\\[5pt]
练习卷：2025-12-1 \qquad 记录日期：2026 年 11 月 1 日
\end{center}

\section{这次练习说明了什么}
第一遍按照考试方式连续听，两题均答错；第二遍不断暂停、回退，并在备忘录中记录听清的内容，最后选择了 B、A。第一遍具体选了什么、重听了多少次，尚未记录。

目前能确认的是：已有不少声音识别和句意理解的基础，但一遍播放时的信息处理还不稳定。听写中已经出现 \en{engine cover}、\en{oil check}、\en{grandson}、\en{breathing} 和 \en{was stuck} 等内容。不能把这次表现概括为“整篇都听不懂”，也不能凭两题判断整体水平。

\textbf{第二遍做对不等于已经能在考试条件下做对。}尤其是第二题，听写仍把 \en{reunited} 留空；答案可能来自直接听懂，也可能来自上下文和选项推断，当前不能替本人确定。

\section{两道题的证据与干扰}
\subsection{第 1 题：B，发现的位置}
选项：\en{Under the engine cover of a man's car.}

截图字幕中的证据是 \en{The black cat was found curled up under the engine cover of David King's car.} 猫是在 David King 的汽车引擎盖下面被发现的。

\begin{itemize}
\item \en{school in Kent}：送孙子上学的背景，不能当成发现猫的具体位置。
\item \en{Mr. King's neighbour}：猫主人的身份；不是说发现猫的车属于邻居。
\item 慈善机构：后续参与救援，不是发现猫的地点。
\end{itemize}

\subsection{第 2 题：A，最终的去向}
选项：\en{It got reunited with its owner.}

截图字幕中的证据是 \en{the four-year-old cat was later reunited with its owner, Mr. King's neighbor.} 猫最终和主人团聚，主人就是 King 先生的邻居。

\textbf{救援者与最终照料者要分开。}由慈善机构救出，不代表后来交给该机构照管；被困、害怕，也不能直接推出救援中受伤。

\subsection{把新闻压缩成事件关系}
送孙子上学后检查机油 $\rightarrow$ 发现引擎盖下的猫 $\rightarrow$ 猫往里面退并被卡住 $\rightarrow$ 个人无法取出 $\rightarrow$ 慈善机构救援 $\rightarrow$ 与主人团聚。

{\small\color{gray}材料边界：题目与作文要求核对自用户提供的原卷第 1 页；听写来自用户备忘录截图，文字对照来自用户提供的字幕截图。本次未获得音频，未独立验证字幕，也未核对录音中的两道问题。答案依据现有材料整理，不冒充官方答案页。}

\newpage
\section{听写差异：哪些影响意思}
以下保留听写中的局部原样，不把不确定的声音擅自补成“听到了”。

\renewcommand{\arraystretch}{1.3}
\begin{tabularx}{\textwidth}{@{}>{\raggedright\arraybackslash}p{0.25\textwidth}>{\raggedright\arraybackslash}p{0.29\textwidth}X@{}}
\toprule
我的记录 & 截图字幕 & 复盘重点\\
\midrule
covered up & curled up & “蜷缩着”，不是“被盖住”；整个词组需重新建立声音与意思的联系。\\
dumping his grandson off & dropping his grandson off & 学习 drop somebody off，表示开车送某人到某处并让其下车。\\
for men to have a look & for me to come have a look & 人称和词的边界未切准；字幕意思是叫“我”过来看。\\
sure was alive & sure it was alive & 漏了 it，主要意思仍可保留；与漏掉关键动作区分。\\
It was even further inter the engine & It reversed even further into the engine & 漏了 reversed，失去“往里面退”的动作；into 表示进入。\\
I try to reach & I tried to reach it & 补足过去时与宾语；动作是尝试够到猫。\\
wasen't have any way i can get it out & there wasn't any way I could get it out & 句子骨架是“没有办法把它弄出来”；注意 there wasn't 与 could。\\
was later *** with his owner & was later reunited with its owner & reunited 是第 2 题关键；its 指猫，不能用是否听出人名替代这项理解。\\
\bottomrule
\end{tabularx}

\subsection{不必把所有差异当成同样严重}
第一句已与截图字幕基本一致。人名 David King、地名 Kent 暂时听不清，可以先用“某人”“某地”占位，继续跟后面的内容。漏掉小词、错认词组、漏掉关键谓语，要分别处理。

目前只能说这些地方发生了识别差异。\textbf{没有音频，不能断定是哪一个音被弱读、连读或省略，也不能确认字幕一定正确。}

\subsection{按意思分块，重新听关键句}
斜线只是练习时的意义分组，不表示录音一定在这里停顿。

\en{The black cat was found / curled up / under the engine cover / of David King's car.}

黑猫被发现／蜷缩着／在引擎盖下面／那辆车属于 David King。

\en{Following a rescue / by UK charity / ... / the four-year-old cat / was later reunited with its owner.}

获救之后／救援方是英国慈善机构／……／这只四岁的猫／后来与主人团聚。

\newpage
\section{下一次怎样练}
\begin{enumerate}
\item \textbf{先做一遍。}正常速度、不暂停；记下选择和把握程度。听前比较选项差异：第 1 题都是地点，第 2 题都是猫的结果。
\item \textbf{只回听卡住的片段。}先不看字幕，每处尝试两三次；仍不清楚就对照文字，避免不断重复同一个猜测。
\item \textbf{区分两类问题。}看文字也不懂：先学词义和搭配；看着懂、听不出：重点练声音识别。再移开文字重听，检查是否能直接获得意思。
\item \textbf{回到整篇。}正常速度、不暂停，只抓发现位置、救援经过、最终去向；隔天再听，检查理解是否保留。
\end{enumerate}

这一篇已经熟悉题目，复听做对主要说明这篇材料有所掌握。\textbf{能否迁移到考试，还要用一篇没做过的材料检查。}

\subsection{本篇值得积累的表达}
\begin{tabularx}{\textwidth}{@{}lX@{}}
\toprule
表达 & 本文中的意思与用法\\
\midrule
a school run & 开车接送孩子上学或放学的行程\\
curl up / be curled up & 蜷缩起来／处于蜷缩状态\\
drop somebody off & 开车送某人到某处并让其下车\\
reverse further into ... & 倒退着进入更深处；本文指猫往引擎里面退\\
be stuck & 被卡住、受困；不等于已经受伤\\
get it out & 把它弄出来；it 指前文的猫\\
be reunited with ... & 与……重聚；本文是与主人团聚\\
\bottomrule
\end{tabularx}

\section{下一项练习：先独立写作文}
\subsection{原卷题目}
\begin{quote}\small
Suppose the student union of your university is collecting opinions on improving its work for the coming year. You are now to write a response by suggesting how it can better enrich student life. You will have 30 minutes to write the essay. You should write at least 120 words but no more than 180 words.
\end{quote}

\textbf{任务：}学生会征求下一年工作的改进意见，请写一篇回应，建议如何更好地丰富学生生活。\textbf{30 分钟，120--180 词。}

第一稿先自己写，不预先套范文。可以把 30 分钟分成审题和列要点 5 分钟、写作 20 分钟、检查 5 分钟；这是练习安排，不是额外考试规定。

交稿后依次检查：是否回应学生会的征求意见、建议是否具体且说明作用、段落是否连贯、语法和搭配是否准确。修改时先保留原意，再逐句说明改动。最后从本人写过并理解的句子里，提炼可替换内容的表达与结构，并注明适用情境。

\textbf{记录状态：}作文初稿、修改稿和可复用表达已补在后文；听力第一遍所选答案、第二题是直接听懂还是推断，仍待本人补充。
\newpage
\section{作文初稿：先看实际完成了什么}
\textbf{正文 166 词，标题不计入，符合 120--180 词要求。}统计时将 \en{college.So} 视为两个词，带撇号的形式按一个词计。本人表示写作用时不长，但未记录起止时间，不据截图时间推定具体耗时。

初稿回应了学生会征求建议的任务，有开头、主体和结尾。主体能辨认出三项想法：组织运动，帮助学生了解专业与未来工作的联系，以及了解学生的日常困难。主要问题是基础语法、搭配和长句组织；不是没有内容，也不是需要再凑字数。

\subsection{初稿存档}
以下按截图转录，去掉屏幕自动换行，保留可辨认的拼写、大小写、标点和语法错误；不是修改稿。
\begin{quote}\small
Some suggestion to enrich students life

I headed that stuent union is collecting suggention to improve work for the coming year, as a students of our college, I want to make some suggestion and make our school better in next year.

I believe the most important thing about next year is how to make students know why we need to study in college.So I think we can prepare some activities such as running around school and know that we can make our body strong before walk into our future work, or make an activity and let students to know how there major can help to get value In future. I believe that we can also prepare to walk into some students’s life to know how the students experience their life and what they to in a whole day, it can let student union know students’s life and organize more activities to teach students solve problems when they face problems in life.

I hope these suggestion can help students get some ideas.
\end{quote}

\subsection{先处理三类问题}
\begin{enumerate}
\item \textbf{句型准确。}重点是 \en{let somebody do}、\en{teach somebody to do} 和 \en{before doing}；不同动词的结构不能混用。
\item \textbf{一句表达一层意思。}开头和主体末尾都用逗号连接了独立句。建议、作用、下一项建议，可以分别成句。
\item \textbf{回到学生生活。}“明白为什么上大学”范围很大；需要落到学生会能组织的活动，以及学生能得到的实际帮助。
\end{enumerate}

\textbf{检查顺序：}内容与字数 $\rightarrow$ 句子是否完整 $\rightarrow$ 动词后面的结构 $\rightarrow$ 名词单复数、代词和拼写。写得快以后，把剩余时间用于这一轮检查。

\newpage
\section{逐句修改：从自己的话出发}
\subsection{标题和开头}
原题目中的 \en{Some suggestion to enrich students life} 可改为：
\en{Some Suggestions for Enriching Student Life}。这里有多项建议，\en{suggestions} 用复数；\en{student life} 可以整体表示学生生活。

原开头有两项独立表达：“听说学生会征求建议”和“作为学生提出建议”。拆开后是：
\begin{quote}
I have heard that the student union is collecting suggestions for the coming year. As a student at our university, I would like to offer some ideas to enrich student life.
\end{quote}
\begin{itemize}
\item \en{headed} 是 \en{head} 的过去式，不表示“听说”；这里用 \en{heard}，也可用 \en{I hear that ...}。
\item \en{stuent} 改为 \en{student}；\en{suggention} 改为 \en{suggestions}。
\item \en{as a students} 改为 \en{as a student}；\en{some suggestion} 改为 \en{some suggestions}。
\item \en{next year} 作时间状语时，前面通常不加 \en{in}。本稿开头用 \en{for the coming year}，不重复时间。
\end{itemize}

\subsection{主体第 1 句：把大目标改成可执行的建议}
\en{make students know why we need to study in college} 的意思可以理解，但它把题目扩大为“上大学的意义”。修改稿删去这一笼统总起句，让后文三项建议直接承担主体。专业与学习意义的想法保留在第 2 项建议中。

\subsection{主体第 2 句：拆开运动与专业活动}
原句从 \en{So I think ...} 一直写到 \en{In future}，把运动、健康、就业、专业价值放在了一起。先分成两项，再分别写活动和作用：
\begin{itemize}
\item 运动：\en{organize regular running activities on campus}；
作用：\en{help students stay healthy and make new friends}。
\item 专业：\en{invite graduates to share how they use what they learned in their majors at work}；
作用：帮助理解所学的价值，为未来职业作准备。
\end{itemize}
\textbf{语法修改：}\en{let students to know} 应为 \en{let students know}，因为 \en{let + 人 + 动词原形}；\en{there major} 中的 \en{there} 应为 \en{their}。\en{before walk} 若保留原结构，应改为 \en{before walking}。但“走进未来工作”整体也不自然，宜直接说 \en{prepare for their future careers}。

\textbf{内容展开：}“regular”“make new friends”和毕业生分享的具体形式，是根据原意补出的细节，不是对原句逐字纠错。\en{prepare some activities} 并非在所有语境中都错；这里表示“组织活动”，用 \en{organize} 更直接。

\subsection{主体第 3 句及结尾：明确调查的用途}
\en{walk into some students's life} 像是把“走进学生生活”直接翻译过去；可写为 \en{learn about students' daily lives and difficulties}。\en{students} 已以 s 结尾，所有格写成 \en{students'}，不是 \en{students's}。

\en{what they to} 改为 \en{what they do}；\en{teach students solve problems} 应为 \en{teach students to solve problems}。修改稿使用 \en{help students solve everyday problems}，意思相近，结构是 \en{help + 人 + (to) do}。再把“了解生活”具体化成调查，并把逗号连接的两个独立句拆开。

结尾的 \en{these suggestion} 改为 \en{these suggestions}；\en{help students get some ideas} 比较模糊，改为说明建议如何帮助学生会了解需求、丰富校园生活。

\newpage
\section{保留原意后的考场版本}
\textbf{正文 136 词，无需依靠标题补足字数。}三项想法来自初稿，活动形式及部分作用作了适度展开。本稿是修改示例，不是本人原始表达，也不作为官方范文或评分结果。

\begin{quote}
I have heard that the student union is collecting suggestions for the coming year. As a student at our university, I would like to offer some ideas to enrich student life.

First, the student union could organize regular running activities on campus. These activities would help students stay healthy and make new friends. Second, it could invite graduates to share how they use what they learned in their majors at work. This would help students understand the value of their studies and prepare for their future careers. Finally, the student union could conduct surveys to learn about students' daily lives and difficulties. Based on the results, it could organize practical activities to help students solve everyday problems.

I hope these suggestions will help the student union better understand students' needs and make campus life more enjoyable.
\end{quote}

\section{可灵活替换的结构}
模板保留的是句子的功能与语法。方括号表示需要按题意填入的部分，不能原样写进作文。

\begin{enumerate}
\item \textbf{说明身份与目的：}\\
\en{As a [身份], I would like to offer some suggestions on [名词或动名词短语].}\\
适用于确有相应身份的建议类回应；本题可填 \en{student} 与 \en{improving campus life}，但其他题目须重新审题。
\item \textbf{建议加作用：}\\
\en{The [组织] could organize [具体活动]. This would help [对象] [动词原形及其余部分].}\\
本题例：\en{The student union could organize sports events. This would help students stay healthy.} 可换活动与收益，但二者要有因果联系。
\item \textbf{邀请谁做什么：}\\
\en{It could invite [人] to [动词原形及其余部分].}\\
本题可以邀请毕业生分享工作经验。注意 \en{invite somebody to do}，与 \en{let somebody do} 不同。
\item \textbf{先了解需求，再采取行动：}\\
\en{The [组织] could conduct surveys to learn about [需求或困难]. Based on the results, it could [具体措施].}\\
适合征求意见、改善服务的任务；后一句必须说明调查结果怎样影响措施。
\end{enumerate}

\subsection{下一轮只练一个能力：建议加作用}
先不背整篇。选“运动”或“了解学生困难”中的一项，不看修改稿，用自己的话写两句：第一句写学生会能做什么，第二句写学生能得到什么。可以沿用 \en{could organize ... / This would help ...}，重点检查动词结构与活动、作用之间的联系。

\end{document}
